%% ============================================================
%% preamble/commands.tex — the paper's mathematical vocabulary
%%
%% Purpose      : every macro the sections may use. Sections define no
%%                macros of their own.
%% Dependencies : etoolbox, amsmath/amssymb/mathtools (packages.tex).
%%
%% NOTATION CONTRACT.
%% This paper repairs two named gaps inside an existing framework, so a
%% referee should be able to read a display here against the corresponding
%% display in the source papers without translating. The macro names and
%% the rendered symbols below are therefore taken from the arXiv LaTeX
%% sources of
%%
%%   [KMU-I] Kramer-Miller--Upton, arXiv:2110.08656,
%%   [KM-ab] Kramer-Miller,        arXiv:2006.04936,
%%   [KM-exp] Kramer-Miller,       arXiv:1909.06905,
%%
%% and NOT chosen afresh. Where the three sources disagree with each other
%% the divergence is recorded at the macro, together with which convention
%% this paper follows and why.
%%
%% MONOCHROME CONTRACT: every macro here typesets in black (colors.tex).
%% ============================================================

\usepackage{xspace}

% ======================================================================
% Number systems and fields.
% KMU-I/KM-ab/KM-exp all define \F, \Z, \Q, \C, \N, \PP as \mathbb{...};
% \PP (not \P, which is LaTeX's pilcrow) is their macro for projective
% space. Kept verbatim so a pasted display compiles either way.
% ======================================================================
\newcommand{\A}{\mathbb{A}}
\newcommand{\C}{\mathbb{C}}
\newcommand{\F}{\mathbb{F}}
\newcommand{\N}{\mathbb{N}}
\newcommand{\PP}{\mathbb{P}}
\newcommand{\Q}{\mathbb{Q}}
\newcommand{\R}{\mathbb{R}}
\newcommand{\Z}{\mathbb{Z}}
\newcommand{\Fq}{\F_q}
\newcommand{\Fqbar}{\overline{\F}_q}

% Script objects. The sources use \calA ... \calZ throughout; only the
% ones this paper needs are defined.
\newcommand{\calA}{\mathcal{A}}
\newcommand{\calD}{\mathcal{D}}
\newcommand{\calB}{\mathcal{B}}
\newcommand{\calE}{\mathcal{E}}
\newcommand{\calF}{\mathcal{F}}
\newcommand{\calL}{\mathcal{L}}
\newcommand{\calO}{\mathcal{O}}
\newcommand{\calP}{\mathcal{P}}
\newcommand{\calR}{\mathcal{R}}
\newcommand{\calS}{\mathcal{S}}

% ======================================================================
% Operators.
%
% \NP and \HP are KMU-I's own \DeclareMathOperator names, and this paper
% writes them with the SAME subscript convention: the base of the
% normalisation rides on the operator, \NP_q(\rho) and \HP_q(\rho), never
% on the argument. KMU-I also carries \cHP (a "contact" Hodge polygon);
% this paper does not use it.
%
% ARGUMENT CONVENTION, where the sources differ. KMU-I and KM-ab both
% write \NP_q(\rho) for the polygon of L(\rho,s) --- the character, not the
% L-function, is the argument --- while Liu--Wan and Zhu write NP_T(f) and
% NP(f) with the defining datum inside. We follow KMU-I/KM-ab, and write
% L(\rho,s) explicitly only where the L-function itself is the subject.
% ======================================================================
% ORDER RELATION. All three sources write "lies on or above" as \succeq,
% on polygons and on slope multisets alike; so does this paper. A bare \ge
% between two polygons would be this paper's invention.
\DeclareMathOperator{\NP}{NP}
\DeclareMathOperator{\HP}{HP}
\DeclareMathOperator{\Swan}{Swan}
\DeclareMathOperator{\Branch}{Branch}
\DeclareMathOperator{\Tr}{Tr}
\DeclareMathOperator{\Gal}{Gal}
\DeclareMathOperator{\Frob}{Frob}
\DeclareMathOperator{\Spec}{Spec}
\DeclareMathOperator{\PGL}{PGL}
\DeclareMathOperator{\rank}{rank}
\DeclareMathOperator{\arcsinh}{arcsinh}
\newcommand{\digitsum}{s_2}
\newcommand{\id}{\mathrm{id}}
\newcommand{\ext}{\mathrm{ext}}       % KMU-I's \ext, as in \rho_P^\ext
\newcommand{\con}{\mathrm{con}}       % KM-ab's O_R^con

% ======================================================================
% The curve, the character, the ramification data.
%
% \Xbar is KM-ab's smooth compactification of X; S = \Xbar \setminus X is
% the set of points at infinity, taken GEOMETRICALLY (Section 5 shows the
% degree count forces that reading).
%
% SWAN CONDUCTOR -- the three sources disagree on the SYMBOL, so this is
% where a referee's eye will snag. [KMU-I] writes d_P, indexed by P in S,
% with delta_P = d_P/p^{n-1} (NOT /(p-1): the (p-1) appears only in the
% pi-adic printing of the local Hodge polygon, whose slopes are
% k(p-1)/delta_P where the q-adic printing has k/d_P). [KM-ab] writes s_Q
% and has no delta at all, bundling ramification into a datum
% (s_Q, e_Q, eps_Q, omega_Q). [KM-exp] writes d_i, indexed by i, as a pole
% order. We follow [KMU-I]: d_P and delta_P, because the Hodge polygon we
% prove a bound for is the one [KMU-I] Sec. 1.2 states in the d_P form, and
% because the pi_q-adic / q-adic conversion in Sec. 3 turns on
% v_pi(p) = d_P(p-1)/delta_P.
% ======================================================================
\newcommand{\Xbar}{\overline{X}}
\newcommand{\Sinf}{S}
\newcommand{\rhoP}{\rho_P^{\ext}}
\newcommand{\dP}{d_P}
\newcommand{\deltaP}{\delta_P}

% ======================================================================
% The p-adic side.
%
% \piq is KMU-I's \pi_q = \pi^{v_p(q)}; the two valuations v_\pi and v_q it
% induces differ by that factor, and every truncation parameter in the
% contact theory is stated in one of them. This paper writes v_\pi, v_q and
% v_p as \vpi, \vq, \vp and NEVER abbreviates which is meant.
% ======================================================================
\newcommand{\piq}{\pi_q}
\newcommand{\vpi}{v_{\pi}}
\newcommand{\vq}{v_q}
\newcommand{\vp}{v_p}
\newcommand{\vtwo}{v_2}
\newcommand{\Rq}{R_q}
\newcommand{\OL}{\calO_L}

% ======================================================================
% The type-2 local objects.
%
% \Up is KMU-I Definition 4.9's operator (1/p)\sigma^{-1}\circ
% \Tr_{E/\sigma E}; at p = 2 we write \Utwo. KM-ab calls the same map U_p
% and its Frobenius \nu rather than \sigma (its Section 4.2); the
% dictionary is in Section 5.
%
% \Adag is A^\dagger; \Am is the weighted growth module A^{m,*}_{\pi,P};
% \Dmod is KM-ab's coefficientwise D = \prod_n p^{b(n)} t^n \OL, which is
% the module this paper actually works in.
% ======================================================================
% FROBENIUS LETTERS -- read this before editing any display in Sec. 2 or 5.
% In [KMU-I], \sigma is the p-Frobenius and the q-Frobenius is its iterate.
% In [KM-ab] and [KM-exp], \nu is the p-Frobenius and \sigma = \nu^a is the
% q-FROBENIUS. So the same letter names different maps in the two papers a
% display may be compared against. This paper uses [KMU-I]'s convention
% (\sigma = p-Frobenius) throughout Sec. 2, states that it is doing so at
% the point of use, and writes \nu explicitly in Sec. 5's dictionary, which
% is the only place [KM-ab]'s formalism is quoted.
\newcommand{\Up}{U_p}
\newcommand{\Utwo}{U_2}
\newcommand{\Adag}{A^{\dagger}}
%
% THREE THINGS CALLED m, KEPT APART BY SHAPE. [KMU-I]'s growth radius is a
% TUPLE and is set in bold, \mathbf{m} = (m_P); [KM-ab]'s \mathbf{m} is a
% SCALAR COUNT (the number of ramified points), which is why we never use a
% bold m for a count and write |S| instead, as [KMU-I] does; and the tail
% index of the operator's support in Sec. 2 is a plain italic m. The bold
% face is load-bearing, not decoration.
\newcommand{\Am}{\calA^{\mathbf{m},*}_{\pi,P}}
\newcommand{\Amt}{\widetilde{\calA}^{\mathbf{m},*}_{\pi}}
% [KM-ab] Sec. 4.2's coefficientwise local module at a type-2 point; it is
% \mathcal{D} there (a glossary symbol), and it is the module in which this
% paper's weight actually lives.
\newcommand{\Dmod}{\calD}
\newcommand{\Edag}{\calE^{\dagger}}
% [KM-ab]'s semi-local objects, used only in Sec. 5.
\newcommand{\Ocon}{\calO_{\calR}^{\con}}

% The weight, the support map, the defect. \wt is the admissible weight
% a(k); \jm is the minimal support index j'(k); \dfct is the defect d(k).
% These are the three functions the whole local half is about, so they get
% names rather than being spelled out at each use.
\newcommand{\wt}{a}
\newcommand{\jm}{j'}
\newcommand{\dfct}{d}
\newcommand{\muP}{\mu(P)}

% ======================================================================
% Inline semantics.
% ======================================================================
\newcommand{\term}[1]{\emph{#1}}
\newcommand{\keyterm}[1]{\textbf{#1}}
\newcommand{\code}[1]{\texttt{#1}}

% Block quotes sit a touch smaller than body copy. This paper quotes its
% sources verbatim more than most, because the p >= 3 hypothesis and the
% obstruction to removing it are both stated by the authors in words, and
% a paraphrase would be the thing under dispute.
\AtBeginEnvironment{quotation}{\small}

% ======================================================================
% The paper's numbers — one source of truth.
%
% Every figure quoted in the text comes from the self-checking programs in
% the replication package. Update HERE only; sections never hard-code a
% number. The comment on each macro names what produced it.
% ======================================================================

% --- the repaired weight and its rate ---------------------------------
\newcommand{\RateGamma}{1/6}              % exact optimal rate (Theorem, sharpness)
\newcommand{\RateGeneral}{1/(2e)}         % same at general odd tame index e
\newcommand{\SelfLoopCoeff}{2}            % c_{2e,2e}, independent of e
\newcommand{\SelfLoopCoeffTwo}{8}         % c_{4e,4e}, independent of e

% --- the exact-rational sweeps (replication: closed-form vs series solve)
\newcommand{\SweepPairs}{355}             % (k,m) pairs, series solve vs closed form
\newcommand{\SweepValK}{60}               % valuation identity checked to k <= 60
\newcommand{\SweepTailK}{600}             % tail estimate checked to k <= 600
\newcommand{\SweepTailM}{80}              % ... and m <= 80
\newcommand{\SweepTightPairs}{150}        % tight pairs found, all k = 2 mod 4, m = 1
\newcommand{\SweepAdmK}{400}              % admissibility swept to k <= 400
\newcommand{\SweepAdmM}{259}              % ... full support m <= 259
\newcommand{\SweepDefectHundred}{17}      % d(100)
\newcommand{\SweepDefectTwoHundred}{33}   % d(200)
\newcommand{\SweepDefectFourHundred}{67}  % d(400)
